\originalsection{高级引用结果的定位}
下表给调用结果的原讲物理页或主论文定位; 本卷的有限模型推导与完整研究定理区分. 原讲L均指MAS.865J 2006. 论文仅引用链接, 不装入源码包.
\begin{longtable}{p{.28\linewidth}p{.63\linewidth}}
\toprule 调用结果 & 来源与定位\\\midrule\endhead
一般多冲突元素互异与最优下界 & \href{https://ocw.mit.edu/courses/mas-865j-quantum-information-science-spring-2006/0d7aea4e65f89ad0fffee1a9e5451828_lecture_11.pdf}{MAS L11第2--4页；唯一冲突五维证明另在正文}\\
网格带币搜索的成功率与维度复杂度 & \href{https://arxiv.org/abs/quant-ph/0402107}{Ambainis--Kempe--Rivosh，定理1、4，第3节及证明第6、9节}\\
完整容错阈值与扩展单元重叠控制 & \href{https://ocw.mit.edu/courses/mas-865j-quantum-information-science-spring-2006/7566f945d7af6676c18401680ede99ff_lecture_7.pdf}{MAS L7第4--8页；本卷仅给单故障模型及粗级联界}\\
一般绝热界及电路--绝热多项式等价 & \href{https://arxiv.org/abs/quant-ph/0405098}{MAS L10第1--5页；Aharonov等，定理1.1及第3节}\\
Hallgren实周期精度与数论复杂度 & \href{https://ocw.mit.edu/courses/mas-865j-quantum-information-science-spring-2006/3142d6105eea0fb1c2c6323c9c3f64e0_lecture_12.pdf}{MAS L12第6--11页；本卷区分基本单位范数与Pell周期}\\
Kuperberg筛法的整体资源界 & \href{https://arxiv.org/abs/quant-ph/0302112}{原论文定理1.1、3.1和第3节，PDF第1--4页}\\
unique-SVP到二面体陪集的几何归约 & \href{https://arxiv.org/abs/cs/0304005}{Regev定理1.1、推论1.2、引理3.2及第3节，PDF第2、5--12页}\\
Shannon编码和Schumacher完整反向界 & \href{https://ocw.mit.edu/courses/mas-865j-quantum-information-science-spring-2006/2b88645ad92377d5c37500fd856d7f9a_lecture_14.pdf}{MAS L14第7--13页；正文给BSC和典型投影构造}\\
主序必要充分判据 & \href{https://arxiv.org/abs/quant-ph/9811053}{Nielsen定理1及证明，PDF第1--3页}\\
混态形成成本的正规化与渐近连续性 & \href{https://arxiv.org/abs/quant-ph/0008134}{Hayden--Horodecki--Terhal，定理1，第I--II节}\\
量子相对熵数据处理 & \href{https://doi.org/10.1007/BF01609351}{Lindblad，Completely positive maps and entropy inequalities，CMP 40 (1975)，147--151页}\\
HSW打包、温和测量、正规化经典容量 & \href{https://ocw.mit.edu/courses/mas-865j-quantum-information-science-spring-2006/4266220604eb311dbba4869e942bcd67_lecture_15.pdf}{MAS L15第1--7页及L16第1--2页；正文给结构与模型}\\
正规化量子容量与解耦可达性 & \href{https://arxiv.org/abs/quant-ph/0304127}{Devetak，定理5与命题7，PDF第11--15页}\\
可降解信道单字母相干信息 & \href{https://arxiv.org/abs/quant-ph/0311131}{Devetak--Shor，第VII节，PDF第13--14页}\\
纠缠辅助经典容量及父协议资源式 & \href{https://arxiv.org/abs/quant-ph/0308044}{MAS L16第3--4、7--8页；Devetak--Harrow--Winter原论文第1--4页}\\
全量子母协议、状态合并与负条件熵 & \href{https://arxiv.org/abs/quant-ph/0606225}{Abeyesinghe等，第III--IV、VII节，PDF第4--13页；状态合并另见0505062}\\
低维PPT充分性与高维束缚纠缠边界 & \href{https://arxiv.org/abs/quant-ph/9605038}{Horodecki等，定理3及证明，PDF第8--9页；MAS L17第7页}\\
私币SMP下界与量子指纹分离 & \href{https://ocw.mit.edu/courses/mas-865j-quantum-information-science-spring-2006/a3023c8cf793563ba7f31391d5135711_lecture_18.pdf}{MAS L18第5--8页，资源模型见正文}\\
一次性量子签名的完整安全定理 & \href{https://arxiv.org/abs/quant-ph/0105032}{Gottesman--Chuang，第6--7节，PDF第4--6页}\\
两体公共品机制的均衡结果 & \href{https://ocw.mit.edu/courses/mas-865j-quantum-information-science-spring-2006/905ee40bf818ef29fa0beca49657c5ae_lecture_19.pdf}{MAS L19第9--10页；本卷仅推导另述GHZ模型}\\
一般攻击下BB84安全与CSS归约 & \href{https://arxiv.org/abs/quant-ph/0003004}{MAS L20第3--7页；Shor--Preskill原论文第2--4页}\\
\bottomrule\end{longtable}
